Considera la funció $f(x)=\sqrt{x-3}$.

\begin{apartats}

\apartat[1,25]{0,75}
Troba el domini de $f$, el punt on talla l'eix $X$ i una primitiva de $f$.

\begin{solucio}
El domini és $[3,+\infty)$, i $f(x)=0$ només per a $x=3$. Escrivint $f(x)=(x-3)^{1/2}$, una primitiva és
$F(x)=\dfrac{(x-3)^{3/2}}{3/2}=\dfrac23\sqrt{(x-3)^3}$.
\end{solucio}

\begin{tria}{area-arrel}
\itemtria{area}{1}{1,25}
Calcula l'àrea de la regió limitada per la gràfica de $f$, l'eix $X$ i la recta $x=4$. I si la recta
fos $x=7$?

\begin{solucio}
Entre 3 i 4, $f\ge0$: $\displaystyle\int_3^4\sqrt{x-3}\,dx=\left[\frac23\sqrt{(x-3)^3}\right]_3^4=\frac23$
unitats quadrades.\\
Fins a $x=7$: $\dfrac23\sqrt{4^3}=\dfrac23\cdot8=\dfrac{16}3$ unitats quadrades.
\end{solucio}

\itemtria{recta-desconeguda}{1}{1,25}
Troba el valor de $b>3$ perquè l'àrea de la regió limitada per la gràfica de $f$, l'eix $X$ i la recta
$x=b$ sigui de 18 unitats quadrades.

\begin{solucio}
L'àrea és $\displaystyle\int_3^b\sqrt{x-3}\,dx=\frac23(b-3)^{3/2}$. Cal $\frac23(b-3)^{3/2}=18$, és a dir,
$(b-3)^{3/2}=27$. Aleshores $b-3=27^{2/3}=9$, i $b=12$.
\end{solucio}
\end{tria}

\begin{nomesllarg}
\apartat{0,75}
Calcula l'àrea de la regió tancada per la paràbola $y=x^2-9$ i l'eix $X$.

\begin{solucio}
Talla l'eix a $x=\pm3$, i entre aquests valors és negativa:
$\left|\displaystyle\int_{-3}^3(x^2-9)\,dx\right|=\left|\left[\frac{x^3}3-9x\right]_{-3}^3\right|=|-18-18|=36$
unitats quadrades.
\end{solucio}
\end{nomesllarg}

\end{apartats}
